% Generated draft: requires mathematical and source review.
\documentclass[11pt]{article}
\usepackage[margin=25mm]{geometry}
\usepackage{fontspec}
\setmainfont{Latin Modern Roman}
\usepackage{amsmath,amssymb,mathtools,unicode-math}
\setmathfont{Latin Modern Math}
\usepackage{enumitem,needspace,xurl,hyperref}
\setlength{\parindent}{0pt}
\setlength{\parskip}{5pt}
\setlength{\emergencystretch}{3em}
\newcommand{\N}{\mathbb{N}}
\newcommand{\Z}{\mathbb{Z}}
\newcommand{\Q}{\mathbb{Q}}
\newcommand{\R}{\mathbb{R}}
\newcommand{\C}{\mathbb{C}}
\newcommand{\F}{\mathbb{F}}
\newcommand{\A}{\mathbb{A}}
\newcommand{\PP}{\mathbb{P}}
\newcommand{\E}{\mathbb{E}}
\newcommand{\eps}{\varepsilon}
\newcommand{\dd}{\,\mathrm{d}}
\begin{document}
\section*{UnsolvedMath ID: 30005409}
\textbf{Source problem code:} catalogue entry
\subsection{Definitions and mathematical statement}
The supplied question does not specify the class of rings or singularities under consideration. The following definitions give a conventional conditional interpretation, without asserting that its hypotheses exhaust the intended scope. If \(R\) is a left noetherian ring, let \(\operatorname{mod}R\) denote the category of finitely generated left \(R\)-modules, and let \(D^b(\operatorname{mod}R)\) be its bounded derived category. Let \(\operatorname{Perf}(R)\) be the full subcategory of objects isomorphic to bounded complexes of finitely generated projective \(R\)-modules. The singularity category is the Verdier quotient
\[
D_{\mathrm{sg}}(R):=D^b(\operatorname{mod}R)/\operatorname{Perf}(R).
\]
Writing \(K^{-,b}(\operatorname{proj}R)\) for the homotopy category of bounded-above complexes of finitely generated projective modules with bounded cohomology, and \(K^b(\operatorname{proj}R)\) for its bounded-complex subcategory, the conventional equivalent presentation is
\[
D_{\mathrm{sg}}(R)\simeq K^{-,b}(\operatorname{proj}R)/K^b(\operatorname{proj}R).
\]
This interprets the damaged quotient formula in the supplied extraction; it is not a transcription verified against the report.

\textbf{Exploratory question.} What information about a singularity represented by \(R\) is determined by \(D_{\mathrm{sg}}(R)\)? For a specified class of admissible rings, this asks which properties or invariants of the singularities must agree whenever
\[
D_{\mathrm{sg}}(R)\simeq D_{\mathrm{sg}}(S)
\]
as triangulated categories. The admissible class, the intended invariants, and whether additional categorical structure must be preserved remain unspecified. Consequently, the record supplies a research question rather than a single fully quantified reconstruction conjecture.
\subsection{Short English statement}
Which features of a singularity can be recovered from its singularity category, formed by quotienting the bounded derived category by perfect complexes? The source leaves the ring hypotheses, desired features, and categorical structure to be retained unspecified.
\subsection{Sources}
\begin{itemize}
\item Local UnsolvedMath registry, supplied complete source record: permanent ID \(30005409\), source problem code \texttt{catalogue entry}, title \emph{Singularity Categories as Invariants of Singularities}. The record belongs to the Oberwolfach Reports Open Problems collection and gives proposed year 2023; no individual proposer is supplied.
\item Supplied report citation: \emph{Resolutions in Local Algebra and Singularity Theory} (2023), Oberwolfach Reports, pp.\ 327--395, DOI \texttt{10.4171/owr/2023/6}. \url{https://doi.org/10.4171/owr/2023/6}. The extracted original text includes the page marker 349.
\item Supplied literature source: ``Kalck, singularity categories.'' \url{https://ems.press/content/serial-article-files/47000}.
\item Supplied literature source: ``DG singular equivalence and singular locus.'' \url{https://arxiv.org/abs/2403.13637}.
\item Literature assessment dated exactly \texttt{2026-08-22}: supplied dated curation, not independently verified here. It labels the question open and exploratory, reports known characterizations detecting regularity, hypersurface periodicity, and some Gorenstein cases, and states that no general reconstruction theorem is established. The supplied background explicitly requires independent verification of time-sensitive status and cited claims.
\item The registry labels the statement \texttt{corrected\_verified} and states that the official report was checked directly and extraction artifacts removed while preserving the mathematical target. This is supplied verification metadata, not an independent verification performed for this draft.
\end{itemize}
\end{document}
